\documentclass[11pt,a4paper]{ctexart}
\usepackage[margin=2.2cm]{geometry}
\usepackage{xcolor}
\usepackage{enumitem}
\usepackage{amssymb}
\usepackage[normalem]{ulem}
\usepackage{titlesec}
\usepackage{hyperref}
\definecolor{navy}{HTML}{1A365D}
\definecolor{ink}{HTML}{1F2933}
\definecolor{muted}{HTML}{66788A}
\definecolor{soft}{HTML}{EEF4FA}
\definecolor{gold}{HTML}{8A5A12}
\hypersetup{colorlinks=true,linkcolor=navy,urlcolor=navy}
\setlist{nosep}
\titleformat{\section}{\Large\bfseries\color{navy}}{}{0pt}{}
\titleformat{\subsection}{\large\bfseries\color{navy}}{}{0pt}{}
\pagestyle{plain}

\title{\textbf{七下名词专项}\\[4pt]
\large 人教版（新课纲版）七年级下册 · 名词（Nouns）}
\author{初中英语资料站}
\date{2026-11-12}
\begin{document}
\maketitle

\section*{使用说明}
\begin{itemize}[leftmargin=1.4em]
\item 本讲义对应人教版（新课纲版）七年级下册语法要求：名词的可数性与复数形式。
\item 共 3 讲语法、6 级 48 题。
\item 在线版：\url{https://1901151458-prog.github.io/english-materials/g7b-nouns/}
\end{itemize}

\section{可数名词和不可数名词}
名词按其所表示的事物的性质可分为\textbf{可数名词}和\textbf{不可数名词}。可数名词有单数和复数两种形式（an apple--two apples、a bag--some bags）；不可数名词一般只有一种形式，没有复数形式（milk、bread、rice）。
\begin{itemize}[leftmargin=1.4em]
\item \textbf{两可词}：fruit、fish、chicken 既可作可数名词，也可作不可数名词。
\end{itemize}

\section{名词复数的构成}
\begin{center}
\begin{tabular}{p{4.8cm}p{3.4cm}p{6.0cm}}
\hline
\textbf{类别} & \textbf{构成} & \textbf{例}\\
\hline
一般情况 & 词尾加 -s & book--books、bag--bags、key--keys、boy--boys\\
以 s、x、ch、sh 结尾 & 词尾加 -es & class--classes、box--boxes、watch--watches、dish--dishes\\
以辅音字母加 y 结尾 & 变 y 为 i 加 -es & family--families、strawberry--strawberries、party--parties\\
\hline
\end{tabular}
\end{center}
\textbf{不规则复数：}man--men、woman--women、child--children、sheep--sheep。

\section{复数的读音}
\begin{center}
\begin{tabular}{p{5.4cm}p{2.6cm}p{5.2cm}}
\hline
\textbf{前一个音} & \textbf{-s/-es 读} & \textbf{例}\\
\hline
清辅音 & /s/ & books /bʊks/、caps /kæps/\\
浊辅音、元音 & /z/ & bags /bæɡz/、keys /kiːz/、boys /bɔɪz/\\
/s/、/z/、/ʃ/、/tʃ/、/dʒ/ & /ɪz/ & classes、boxes、watches、dishes\\
\hline
\end{tabular}
\end{center}


\newpage
\section{分级练习（48 题）}
\noindent 按难度分级排列，先做基础，再做成分与扩展，最后做综合题。

\subsection*{第 1 级 · 可数与不可数}
\noindent\textcolor{gray}{可数名词有复数；不可数名词一般没有。}\par\smallskip
\begin{enumerate}[leftmargin=*,itemsep=4pt]
\item apple 属于 \_\_\_。
\begin{enumerate}[label=\Alph*.,leftmargin=2em,itemsep=1pt,topsep=2pt]
\item 可数名词
\item 不可数名词
\item 专有名词
\item 抽象名词
\end{enumerate}
\item milk 属于 \_\_\_。
\begin{enumerate}[label=\Alph*.,leftmargin=2em,itemsep=1pt,topsep=2pt]
\item 可数名词
\item 不可数名词
\item 专有名词
\item 复数名词
\end{enumerate}
\item bread 属于 \_\_\_。
\begin{enumerate}[label=\Alph*.,leftmargin=2em,itemsep=1pt,topsep=2pt]
\item 可数名词
\item 不可数名词
\item 专有名词
\item 数词
\end{enumerate}
\item rice 属于 \_\_\_。
\begin{enumerate}[label=\Alph*.,leftmargin=2em,itemsep=1pt,topsep=2pt]
\item 可数名词
\item 不可数名词
\item 专有名词
\item 复数名词
\end{enumerate}
\item a bag–some \_\_\_
\begin{enumerate}[label=\Alph*.,leftmargin=2em,itemsep=1pt,topsep=2pt]
\item bag
\item bags
\item bag’s
\item bagss
\end{enumerate}
\item 不可数名词一般 \_\_\_ 复数形式。
\begin{enumerate}[label=\Alph*.,leftmargin=2em,itemsep=1pt,topsep=2pt]
\item 没有
\item 有
\item 有两种
\item 只有 -s
\end{enumerate}
\item an apple–two \_\_\_
\begin{enumerate}[label=\Alph*.,leftmargin=2em,itemsep=1pt,topsep=2pt]
\item apple
\item apples
\item apple’s
\item applies
\end{enumerate}
\item 可数名词有 \_\_\_ 两种形式。
\begin{enumerate}[label=\Alph*.,leftmargin=2em,itemsep=1pt,topsep=2pt]
\item 单数和复数
\item 主格和宾格
\item 原级和比较级
\item 肯定和否定
\end{enumerate}
\end{enumerate}

\subsection*{第 2 级 · 复数 · 规则变化}
\noindent\textcolor{gray}{加 -s、加 -es、y 变 i 加 -es。}\par\smallskip
\begin{enumerate}[leftmargin=*,itemsep=4pt]
\item book 的复数是 \_\_\_。
\begin{enumerate}[label=\Alph*.,leftmargin=2em,itemsep=1pt,topsep=2pt]
\item books
\item bookes
\item book’s
\item bookies
\end{enumerate}
\item class 的复数是 \_\_\_。
\begin{enumerate}[label=\Alph*.,leftmargin=2em,itemsep=1pt,topsep=2pt]
\item classes
\item classs
\item classies
\item class’es
\end{enumerate}
\item box 的复数是 \_\_\_。
\begin{enumerate}[label=\Alph*.,leftmargin=2em,itemsep=1pt,topsep=2pt]
\item boxes
\item boxs
\item boxies
\item box’es
\end{enumerate}
\item watch 的复数是 \_\_\_。
\begin{enumerate}[label=\Alph*.,leftmargin=2em,itemsep=1pt,topsep=2pt]
\item watches
\item watchs
\item watchies
\item watch’es
\end{enumerate}
\item dish 的复数是 \_\_\_。
\begin{enumerate}[label=\Alph*.,leftmargin=2em,itemsep=1pt,topsep=2pt]
\item dishes
\item dishs
\item dishies
\item dish’es
\end{enumerate}
\item family 的复数是 \_\_\_。
\begin{enumerate}[label=\Alph*.,leftmargin=2em,itemsep=1pt,topsep=2pt]
\item families
\item familys
\item familyes
\item familiees
\end{enumerate}
\item strawberry 的复数是 \_\_\_。
\begin{enumerate}[label=\Alph*.,leftmargin=2em,itemsep=1pt,topsep=2pt]
\item strawberries
\item strawberrys
\item strawberryes
\item strawberries’
\end{enumerate}
\item key 的复数是 \_\_\_。
\begin{enumerate}[label=\Alph*.,leftmargin=2em,itemsep=1pt,topsep=2pt]
\item keys
\item keies
\item keyes
\item key’s
\end{enumerate}
\end{enumerate}

\subsection*{第 3 级 · 复数 · 读音}
\noindent\textcolor{gray}{/s/、/z/、/ɪz/ 的判断。}\par\smallskip
\begin{enumerate}[leftmargin=*,itemsep=4pt]
\item books 中 -s 读 \_\_\_。
\begin{enumerate}[label=\Alph*.,leftmargin=2em,itemsep=1pt,topsep=2pt]
\item /s/
\item /z/
\item /ɪz/
\item /t/
\end{enumerate}
\item bags 中 -s 读 \_\_\_。
\begin{enumerate}[label=\Alph*.,leftmargin=2em,itemsep=1pt,topsep=2pt]
\item /s/
\item /z/
\item /ɪz/
\item /d/
\end{enumerate}
\item classes 中 -es 读 \_\_\_。
\begin{enumerate}[label=\Alph*.,leftmargin=2em,itemsep=1pt,topsep=2pt]
\item /s/
\item /z/
\item /ɪz/
\item /es/
\end{enumerate}
\item boxes 中 -es 读 \_\_\_。
\begin{enumerate}[label=\Alph*.,leftmargin=2em,itemsep=1pt,topsep=2pt]
\item /s/
\item /z/
\item /ɪz/
\item /əz/
\end{enumerate}
\item keys 中 -s 读 \_\_\_。
\begin{enumerate}[label=\Alph*.,leftmargin=2em,itemsep=1pt,topsep=2pt]
\item /s/
\item /z/
\item /ɪz/
\item /iːz/
\end{enumerate}
\item watches 中 -es 读 \_\_\_。
\begin{enumerate}[label=\Alph*.,leftmargin=2em,itemsep=1pt,topsep=2pt]
\item /s/
\item /z/
\item /ɪz/
\item /tʃɪz/
\end{enumerate}
\item boys 中 -s 读 \_\_\_。
\begin{enumerate}[label=\Alph*.,leftmargin=2em,itemsep=1pt,topsep=2pt]
\item /s/
\item /z/
\item /ɪz/
\item /ɔɪz/
\end{enumerate}
\item dishes 中 -es 读 \_\_\_。
\begin{enumerate}[label=\Alph*.,leftmargin=2em,itemsep=1pt,topsep=2pt]
\item /s/
\item /z/
\item /ɪz/
\item /ʃz/
\end{enumerate}
\end{enumerate}

\subsection*{第 4 级 · 复数 · 不规则}
\noindent\textcolor{gray}{man-men、child-children、sheep-sheep 等。}\par\smallskip
\begin{enumerate}[leftmargin=*,itemsep=4pt]
\item man 的复数是 \_\_\_。
\begin{enumerate}[label=\Alph*.,leftmargin=2em,itemsep=1pt,topsep=2pt]
\item men
\item mans
\item manes
\item mens
\end{enumerate}
\item woman 的复数是 \_\_\_。
\begin{enumerate}[label=\Alph*.,leftmargin=2em,itemsep=1pt,topsep=2pt]
\item women
\item womans
\item womens
\item womies
\end{enumerate}
\item child 的复数是 \_\_\_。
\begin{enumerate}[label=\Alph*.,leftmargin=2em,itemsep=1pt,topsep=2pt]
\item children
\item childs
\item childes
\item childrens
\end{enumerate}
\item sheep 的复数是 \_\_\_。
\begin{enumerate}[label=\Alph*.,leftmargin=2em,itemsep=1pt,topsep=2pt]
\item sheep
\item sheeps
\item sheepes
\item sheep’s
\end{enumerate}
\item 下列哪个是不规则复数？
\begin{enumerate}[label=\Alph*.,leftmargin=2em,itemsep=1pt,topsep=2pt]
\item man–men
\item book–books
\item bag–bags
\item key–keys
\end{enumerate}
\item child 的复数属于 \_\_\_ 变化。
\begin{enumerate}[label=\Alph*.,leftmargin=2em,itemsep=1pt,topsep=2pt]
\item 不规则
\item 规则
\item 加 -s
\item 加 -es
\end{enumerate}
\item two \_\_\_（woman）
\begin{enumerate}[label=\Alph*.,leftmargin=2em,itemsep=1pt,topsep=2pt]
\item women
\item womans
\item womens
\item woman
\end{enumerate}
\item three \_\_\_（sheep）
\begin{enumerate}[label=\Alph*.,leftmargin=2em,itemsep=1pt,topsep=2pt]
\item sheep
\item sheeps
\item sheepes
\item sheeps’
\end{enumerate}
\end{enumerate}

\subsection*{第 5 级 · 既可数又不可数}
\noindent\textcolor{gray}{fruit、fish、chicken 等。}\par\smallskip
\begin{enumerate}[leftmargin=*,itemsep=4pt]
\item fruit 可以 \_\_\_。
\begin{enumerate}[label=\Alph*.,leftmargin=2em,itemsep=1pt,topsep=2pt]
\item 既作可数也作不可数
\item 只能作可数
\item 只能作不可数
\item 不能作名词
\end{enumerate}
\item fish 可以 \_\_\_。
\begin{enumerate}[label=\Alph*.,leftmargin=2em,itemsep=1pt,topsep=2pt]
\item 既作可数也作不可数
\item 只能作可数
\item 只能作不可数
\item 只能作动词
\end{enumerate}
\item chicken 可以 \_\_\_。
\begin{enumerate}[label=\Alph*.,leftmargin=2em,itemsep=1pt,topsep=2pt]
\item 既作可数也作不可数
\item 只能作可数
\item 只能作不可数
\item 不能作名词
\end{enumerate}
\item 下列哪个词既可数又不可数？
\begin{enumerate}[label=\Alph*.,leftmargin=2em,itemsep=1pt,topsep=2pt]
\item fish
\item book
\item milk
\item rice
\end{enumerate}
\item 下列哪个是不可数名词？
\begin{enumerate}[label=\Alph*.,leftmargin=2em,itemsep=1pt,topsep=2pt]
\item rice
\item apple
\item bag
\item pen
\end{enumerate}
\item 下列哪个是可数名词？
\begin{enumerate}[label=\Alph*.,leftmargin=2em,itemsep=1pt,topsep=2pt]
\item apple
\item milk
\item bread
\item rice
\end{enumerate}
\item some \_\_\_（bag）
\begin{enumerate}[label=\Alph*.,leftmargin=2em,itemsep=1pt,topsep=2pt]
\item bags
\item bag
\item bag’s
\item bagies
\end{enumerate}
\item two \_\_\_（apple）
\begin{enumerate}[label=\Alph*.,leftmargin=2em,itemsep=1pt,topsep=2pt]
\item apples
\item apple
\item apple’s
\item applies
\end{enumerate}
\end{enumerate}

\subsection*{第 6 级 · 综合提升}
\noindent\textcolor{gray}{可数性、复数形式与读音混合辨析。}\par\smallskip
\begin{enumerate}[leftmargin=*,itemsep=4pt]
\item 名词按其性质分为 \_\_\_ 和不可数名词。
\begin{enumerate}[label=\Alph*.,leftmargin=2em,itemsep=1pt,topsep=2pt]
\item 可数名词
\item 专有名词
\item 普通名词
\item 抽象名词
\end{enumerate}
\item class 的复数 \_\_\_。
\begin{enumerate}[label=\Alph*.,leftmargin=2em,itemsep=1pt,topsep=2pt]
\item classes
\item classs
\item classies
\item classees
\end{enumerate}
\item party 的复数 \_\_\_。
\begin{enumerate}[label=\Alph*.,leftmargin=2em,itemsep=1pt,topsep=2pt]
\item parties
\item partys
\item partyies
\item party’s
\end{enumerate}
\item book–books 中 -s 读 \_\_\_。
\begin{enumerate}[label=\Alph*.,leftmargin=2em,itemsep=1pt,topsep=2pt]
\item /s/
\item /z/
\item /ɪz/
\item /b/
\end{enumerate}
\item box–boxes 中 -es 读 \_\_\_。
\begin{enumerate}[label=\Alph*.,leftmargin=2em,itemsep=1pt,topsep=2pt]
\item /ɪz/
\item /s/
\item /z/
\item /es/
\end{enumerate}
\item two \_\_\_（child）
\begin{enumerate}[label=\Alph*.,leftmargin=2em,itemsep=1pt,topsep=2pt]
\item children
\item childs
\item childes
\item childrens
\end{enumerate}
\item milk 有复数形式吗？
\begin{enumerate}[label=\Alph*.,leftmargin=2em,itemsep=1pt,topsep=2pt]
\item 没有
\item 有
\item 有两种
\item 只有 milk’s
\end{enumerate}
\item family–families 中 -ies 读 \_\_\_。
\begin{enumerate}[label=\Alph*.,leftmargin=2em,itemsep=1pt,topsep=2pt]
\item /ɪz/
\item /s/
\item /z/
\item /aɪz/
\end{enumerate}
\end{enumerate}

\newpage
\section*{答案与解析}
\small
\begin{enumerate}[leftmargin=1.6em,itemsep=2pt]
\item \textbf{Q1-1 A}：apple 有复数 apples；教材：可数名词有单数和复数两种形式。
\item \textbf{Q1-2 B}：milk 一般只有一种形式，不可数；教材例句。
\item \textbf{Q1-3 B}：bread 不可数；教材例句。
\item \textbf{Q1-4 B}：rice 不可数；教材例句。
\item \textbf{Q1-5 B}：可数名词有复数形式；教材例句 a bag–some bags。
\item \textbf{Q1-6 A}：教材：不可数名词一般只有一种形式，没有复数形式。
\item \textbf{Q1-7 B}：可数名词变复数；教材例句 an apple–two apples。
\item \textbf{Q1-8 A}：教材：可数名词有单数和复数两种形式。
\item \textbf{Q2-1 A}：一般情况词尾加 -s；教材表1。
\item \textbf{Q2-2 A}：以 s 结尾加 -es；教材表1。
\item \textbf{Q2-3 A}：以 x 结尾加 -es；教材表1。
\item \textbf{Q2-4 A}：以 ch 结尾加 -es；教材表1。
\item \textbf{Q2-5 A}：以 sh 结尾加 -es；教材表1。
\item \textbf{Q2-6 A}：辅音字母 + y 结尾，y 变 i 加 -es；教材表1。
\item \textbf{Q2-7 A}：y 变 i 加 -es；教材表1。
\item \textbf{Q2-8 A}：元音字母 + y 结尾，直接加 -s；教材表1 keys。
\item \textbf{Q3-1 A}：清辅音 /k/ 后读 /s/；教材表1。
\item \textbf{Q3-2 B}：浊辅音 /ɡ/ 后读 /z/；教材表1。
\item \textbf{Q3-3 C}：在 /s/ 音后读 /ɪz/；教材表1。
\item \textbf{Q3-4 C}：在 /s/ 音后读 /ɪz/；教材表1。
\item \textbf{Q3-5 B}：元音后读 /z/；教材表1。
\item \textbf{Q3-6 C}：在 /tʃ/ 后读 /ɪz/；教材表1。
\item \textbf{Q3-7 B}：元音后读 /z/；教材表1。
\item \textbf{Q3-8 C}：在 /ʃ/ 后读 /ɪz/；教材表1。
\item \textbf{Q4-1 A}：不规则复数；教材注释。
\item \textbf{Q4-2 A}：不规则复数；教材注释。
\item \textbf{Q4-3 A}：不规则复数；教材注释。
\item \textbf{Q4-4 A}：单复数同形；教材注释。
\item \textbf{Q4-5 A}：教材注释列出 man–men、woman–women、child–children、sheep–sheep。
\item \textbf{Q4-6 A}：child–children 属不规则变化。
\item \textbf{Q4-7 A}：woman 的复数是不规则的 women。
\item \textbf{Q4-8 A}：sheep 单复数同形。
\item \textbf{Q5-1 A}：教材：fruit 既可充当可数名词也可充当不可数名词。
\item \textbf{Q5-2 A}：教材：fish 两可；同时也可作动词。
\item \textbf{Q5-3 A}：教材：chicken 两可（鸡／鸡肉）。
\item \textbf{Q5-4 A}：教材列出 fruit、fish、chicken 三个两可词。
\item \textbf{Q5-5 A}：rice 不可数；教材例句。
\item \textbf{Q5-6 A}：apple 可数；教材例句。
\item \textbf{Q5-7 A}：可数名词复数为 bags；教材例句 some bags。
\item \textbf{Q5-8 A}：可数名词变复数；教材例句 two apples。
\item \textbf{Q6-1 A}：教材：名词按其所表示的事物的性质可分为可数名词和不可数名词。
\item \textbf{Q6-2 A}：以 s 结尾加 -es。
\item \textbf{Q6-3 A}：辅音字母 + y，y 变 i 加 -es。
\item \textbf{Q6-4 A}：清辅音 /k/ 后读 /s/。
\item \textbf{Q6-5 A}：在 /s/ 音后读 /ɪz/。
\item \textbf{Q6-6 A}：不规则复数 children。
\item \textbf{Q6-7 A}：不可数名词一般没有复数形式。
\item \textbf{Q6-8 A}：在 /z/ 等音后读 /ɪz/；教材表1 /ˈfæməlɪz/。
\end{enumerate}

\end{document}
